% Generated draft: requires mathematical and source review.
\documentclass[11pt]{article}
\usepackage[margin=25mm]{geometry}
\usepackage{fontspec}
\setmainfont{Latin Modern Roman}
\usepackage{amsmath,amssymb,mathtools,unicode-math}
\setmathfont{Latin Modern Math}
\usepackage{enumitem,needspace,xurl,hyperref}
\setlength{\parindent}{0pt}
\setlength{\parskip}{5pt}
\setlength{\emergencystretch}{3em}
\newcommand{\N}{\mathbb{N}}
\newcommand{\Z}{\mathbb{Z}}
\newcommand{\Q}{\mathbb{Q}}
\newcommand{\R}{\mathbb{R}}
\newcommand{\C}{\mathbb{C}}
\newcommand{\F}{\mathbb{F}}
\newcommand{\A}{\mathbb{A}}
\newcommand{\PP}{\mathbb{P}}
\newcommand{\E}{\mathbb{E}}
\newcommand{\eps}{\varepsilon}
\newcommand{\dd}{\,\mathrm{d}}
\begin{document}
\section*{UnsolvedMath ID: 30003055}
\textbf{Source problem code:} catalogue entry
\subsection{Definitions and mathematical statement}
The supplied record describes a collection of comparative research questions, rather than a single quantified conjecture. Its mathematical setting is not sufficiently specified to determine a unique theorem statement. The following terminology and conditional notation organize the questions without selecting missing hypotheses.

A polynomial spline is a function that is polynomial on each cell of a specified partition and satisfies specified continuity conditions across cell interfaces. Maximal and minimal continuity must be understood relative to the polynomial degree, partition, and admissible continuity requirements; these data are not given. A finite-element approximation space likewise requires a domain, a mesh, local approximation spaces, and interface conditions. The record does not specify which classical finite elements are intended.

For an open set \(\Omega\subset\R^d\), the divergence space is
\[
H(\operatorname{div};\Omega)=\{v\in L^2(\Omega;\R^d):\operatorname{div}v\in L^2(\Omega)\},
\]
where divergence is interpreted distributionally. A divergence-conforming spline space is a spline space contained in this space. In the usual three-dimensional vector formulation,
\[
H(\operatorname{curl};\Omega)=\{v\in L^2(\Omega;\R^3):\operatorname{curl}v\in L^2(\Omega;\R^3)\},
\qquad \Omega\subset\R^3,
\]
and a curl-conforming spline space is a spline space contained in this space. This definition explains the terminology only: it does not restrict the source question to dimension three. Curl conventions in other dimensions and any boundary conditions remain unspecified.

For a specified normed function space \(X\), a target \(u\in X\), and a finite-dimensional approximation space \(V\subset X\), one possible accuracy quantity is the best-approximation error
\[
E_X(u,V):=\inf_{v\in V}\|u-v\|_X.
\]
The source does not choose this quantity over the error of a particular numerical method. It also does not define robustness, an error norm, a target regularity class, or a common computational budget. A space and a basis for that space are distinct objects: changing a basis leaves \(E_X(u,V)\) unchanged but can affect numerical conditioning and computational cost.

\textbf{Research questions retained from the source.} Determine how accurate and robust maximally continuous, minimally continuous, divergence-conforming, and curl-conforming spline approximations are; compare spline approximations with classical finite-element approximations; determine how these comparisons depend on dimension; and determine which polynomial and Fourier bases are preferable for multidimensional approximation and when those preferences admit quantitative justification. Tensor-product settings are also mentioned in the supplied literature assessment.

These requests require a specified approximation problem and comparison criterion before they become precise mathematical assertions. In particular, no universal ordering of the listed spaces or bases is asserted here. The intended domains, dimensions, polynomial degrees, meshes, target classes, boundary or periodicity conditions, algorithms, resource constraints, and robustness parameters are unresolved.
\subsection{Short English statement}
This item asks for quantitative comparisons of spline and classical finite-element approximation, including continuity, divergence and curl conformity, dimension, and polynomial versus Fourier bases. It is a broad research agenda: the supplied record does not specify the mathematical setting or criteria needed for a single precise conjecture. Its supplied open status requires independent review.
\subsection{Sources}
\begin{itemize}
\item Local UnsolvedMath registry: permanent ID \(30003055\), title \emph{Comparative Accuracy of Spline and Finite-Element Bases}, source code \texttt{catalogue entry}, in the Oberwolfach Reports Open Problems collection. The supplied record gives proposed year 2016 and names no individual proposer.
\item Supplied report citation: \emph{Mini-Workshop: Mathematical Foundations of Isogeometric Analysis} (2016), Oberwolfach Reports, pp.~341--385, DOI \texttt{10.4171/owr/2016/8}. \url{https://doi.org/10.4171/owr/2016/8}. The supplied background describes these reports as documentation of mathematical workshops and their open-problem discussions; it supplies no item-specific page or speaker attribution.
\item Supplied literature source, labeled ``original OWR report'': \url{https://doi.org/10.4171/owr/2016/8}.
\item Supplied literature source, labeled ``recent spline-basis comparison'': \url{https://doi.org/10.1016/j.cma.2023.116659}. No authors, article title, or specific supporting result are supplied.
\item Supplied literature assessment dated exactly \texttt{2026-08-21}: the item is labeled open and described as a broad collection of unresolved numerical-analysis comparisons across spline, finite-element, divergence, curl, tensor-product, and Fourier settings, with no single theorem answering the full list. This is supplied dated curation, not an independently verified literature assessment. The background explicitly requests independent verification of time-sensitive status and cited claims.
\end{itemize}
\end{document}
